\section{Observations} \label{sec:observations}

\eric{brief summary of observations: commissioning, first alerts, ongoing data taking.  
What we expect to be in the PPDB.}

Unlike Rubin's Data Releases (e.g., DP1, DP2 \eric{cite}), which contain a fixed dataset acquired before the processing begins, Alert Production provides a continuous, ongoing release of data obtained by Rubin that is released within seconds to days of observation.
During commissioning, alerts were transmitted to community alert brokers for technical integration tests beginning on \AlertIntegrationStartDate.
Since data from this period are available in alert history and the Prompt Products Database (PPDB; \S \ref{sec:ppdb}) \eric{image release?}, we include them in this paper, although they were not available publicly at the time of processing.
Real-time public alert release began on \PublicAlertsStartDate.
In this paper we include data obtained until \eric{final date}.

During the commissioning and early operations phases, scientific data taking was intermixed with ongoing technical commisisoning activities that did not produce data suitable for alert production.
Calibration, delivered image quality, and pipeline processing were under substantial development.
Additionally, template images were only available initially for a small subset of the sky area.
Accordingly, the quantity and quality of early alerts produced are not representive of steady-state performance, and not all data products are available.
We specify the range of data included in our performance analysis (\S \ref{sec:performance}).




\eric{consider a figure or histogram showing the alert rate over time, perhaps a sky plot of alert density?? or colored by earliest alert date?}

\eric{scheduler discussion/citation?  link to survey-strategy.lsst.io?  mention status dashboards at prompt-products.lsst.io?}